\SourceLeaf{530}beigegeben war und die eine Berichtigung des Art. 139 Disqu. Arithm. enthielt, ist dort der betreffenden Stelle eingefügt. Die Note auf Seite 91 ist einer handschriftlichen Notiz entlehnt. Ausserdem unterscheidet sich die vorliegende Ausgabe von den früheren nur durch die Berichtigung einiger Druckfehler. Die von mir hinzugefügten Einschaltungen sind durch eckige Klammern [\,] kenntlich gemacht.
\par Die \textit{Tafel des quadratischen Characters der Primzahlen} ist nach der Weise der in Art. 99 beschriebenen und (in Art. 331) zur Zerlegung der Zahlen vorzugsweise angewandten Tabula II der Disqu. Arithm. gedruckt. Die Handschrift unter dem Titel ‘\textit{Quadratorum numeris primis divisorum residua lateralia}’ hat in den Schriftzügen am meisten Aehnlichkeit mit der des zweiten Theiles der Tafel zur Verwandlung gemeiner Brüche in Decimalbrüche, sie enthält an der Stelle der den Quadratischen Rest anzeigenden horizontalen Striche kleine Kreise, von denen immer diejenigen durch Linien verbunden sind, die in benachbarten horizontalen oder verticalen Reihen vorkommen. Bei der Correctur wurde ich auf mehrere Fehler aufmerksam, habe dann bei einer einmaligen Vergleichung mit \textsc{Jacobi}’s \textit{Canon Arithmeticus} 190 Abweichungen in den Angaben der Charactere und nach directer Bestimmung diese in Uebereinstimmung mit jenen gedruckten Tafeln gefunden, dem entsprechend ist hier die Ausgabe berichtigt.
\par Von der \textit{Tafel zur Verwandlung gemeiner Brüche in Decimalbrüche} ist hier der erste Theil der Tabula III. der Disqu. Arithm. ähnlich eingerichtet, er enthält für die Primzahlen und deren Potenzen $p^\pi$, welche zwischen 3 und 463 liegen, die Mantissen $(1)$, $(2)\dotstwo(0)$ der Decimalbrüche von $\frac{10\pdot r}{p^\pi}$, $\frac{10\pdot rr}{p^\pi}\dotstwo\frac{10}{p^\pi}$, worin $r$ die Einheit bedeutet, also $(1)=(2)=\dotstwo(0)$ wird, wenn 10 Primitivwurzel von $p^\pi$ ist, sonst aber $r$ die kleinste unter denjenigen Primitivwurzeln von $p^\pi$ bezeichnet, für welche als Basis der Index von 10 den kleinsten Werth annimmt. Die von 1 verschiedenen Werthe von $r$ habe ich zur Erleichterung des Gebrauchs auf Seite 420 der Tafel beigefügt. Die Handschrift, in der auch noch nicht die Unterscheidungsziffern der verschiedenen Perioden angegeben sind, entspricht äusserlich am meisten der Analysis residuorum und scheint in der Zeit dem hier als zweiten Theil der ganzen Tafel hingestellten Stücke voraufzugehen.
Dieser zweite Theil enthält für die Primzahlen und deren Potenz $p^\pi$ zwischen 467 und 997 die Mantissen der Decimalbrüche von $\frac{100}{p^\pi}$. Die Handschrift gibt die Theiler in abnehmender Reihenfolge und schliesst mit den Worten: \textit{Explicitus October} 11. 1795. Im Drucke habe ich beim Theiler 191 Periode $(1)$ die 71\textsuperscript{ste} Ziffer hinzugefügt und beim Theiler 829 eine zwischen der 151 und 152\textsuperscript{sten} Ziffer stehende Zahl fortgelassen.
\par Die von \textsc{Gauss} selbst in einem Briefe (Seite 444) erläuterte \textit{Tafel der Frequenz der Primzahlen} besteht für ihren ersten Theil, welche die Anzahl der Primzahlen in jedem der 1000 ersten Chiliaden gibt in einer Handschrift von \textsc{Gauss}, es finden sich im Nachlass aber nicht die in dem Briefe angedeuteten Abzählungen der der ersten Million angehörenden Hunderte, die eine bestimmte Anzahl von Primzahlen enthalten. Der andere Theil der Tafel nemlich für die zweite und dritte Million ist einer von \textsc{Goldschmidt} allein herrührenden Handschrift entlehnt. Herr \textsc{Meissel} hat durch Abzählung und durch seine Formel die folgenden Berichtigungen zu Seite 436 und 437 gefunden:
